\chapter{The Stealth Codex---Flowing Privacy}

% Reading progression for this chapter:
% Curse (problem) -> Arts (three forms) -> Social consequences of the arts (trust/property) -> Cost (compute/entropy) -> New dimension (time) -> Dao (philosophy) -> Humanity (love)

\begin{myquote}
\textit{"Know the white, yet keep to the black---become the pattern for all under heaven."}

\hfill\textit{---Tao Te Ching}
\end{myquote}

\section{Grandet Sitting on a Mountain of Gold---Data That Cannot Flow}

The life of data lies in its flow.

The velocity of money determines the prosperity of an economy. If the data each of us possesses could flow freely, it too should create value---and the faster the flow, the greater the value. Yet in the Fifth Dimension, there exists a paradox of data value: \textbf{the more valuable the data, the less it dares to flow.}

\vspace{1em}

In the physical world, transactions follow the principle of clean exchange---goods for payment, with property rights clearly defined.
The shadow-clone arts of the Fifth Dimension have shattered those rules entirely.
Once data is handed over, the recipient can copy it ten thousand times at will, and the sender has no way of knowing.
\textbf{If copying is possessing, then establishing ownership loses all meaning.}

\vspace{1em}

The dilemma is this: the value of data is often proportional to its sensitivity.

The most valuable medical data can decode the secrets of life, yet it is also the softest underbelly of privacy;

the most valuable financial data can reveal the pulse of the market, yet it is also the most closely guarded commercial secret.

And so we fall into the trap of Grandet:
data hoarded to death cannot forge connections, and data without connections becomes a lifeless thing.
Data becomes Schr\"{o}dinger's asset---without opening the box, it is simultaneously a potential gold mine and a pile of useless scrap.

\vspace{1em}

Until the Stealth Codex appeared and forged a lock for our data.

This chapter tells the story of privacy-preserving technologies. It does not teach us how the lock is made, but rather helps us see: once these locks exist, how will the balance of power between individuals and platforms, persons and organizations, data and property rights be transformed?


\section{The Invisibility Art of the Fifth Dimension---Data Unseen Yet Free to Flow}

We need a new kind of sorcery.

It must be able to prove truth or falsehood without exposing the data itself, perform computation while keeping data invisible, and establish consensus among parties who do not trust one another.

\vspace{1em}

Its core mantra is just six characters:

\begin{importantbox}
\textbf{Usable yet invisible.}
\end{importantbox}


Legend has it that the art of stealth concealment could render its practitioners invisible between heaven and earth, traveling freely across the boundary of yin and yang, completing their mission without being detected. In the Fifth Dimension, this is precisely the ability we need---to let data flow without being seen.

\vspace{1em}

The Stealth Codex comprises three forms:

The First Form is the Art of Verification, solving the problem of how to prove a capability without revealing it;

The Second Form is the Art of Computation, enabling data to be used through calculation without being understood;

The Third Form is the Art of the Macro, able to access all data yet revealing only the whole, never the individual.

\vspace{1em}

What the three forms share is not algorithmic detail, but a philosophy of separation:
using mathematics to decouple rights and risks that were once bound together.
When the preimage is invisible, value flows more freely.

\vspace{1em}

Let us study each in turn.

\section{First Form of the Stealth Codex: The Art of Verification}

Here lies an irreconcilable contradiction: we need data to know you, yet we must also protect your privacy.

The central proposition of the Art of Verification is: \textbf{prove it is you without revealing your information.}

This is the invisibility art of identity and eligibility.

It allows you to prove ``I am qualified'' without exposing ``who I am.''

\newpage
\subsection{Ali Baba's Cave: Proving You Know the Spell Without Uttering a Word}

Imagine a circular cave with a magic door in the middle that can only be opened by a secret spell.
The prover, Peggy, wants to prove to the verifier, Victor, that she knows the spell---but she refuses to say it aloud.

Perhaps the spell loses its power the moment it is spoken; perhaps Victor does not trust the verifier;

perhaps the spell itself is Victor's most precious secret.

\vspace{1em}
The solution is surprisingly simple.

The verifier waits at the entrance while the prover randomly enters the depths of the cave from either the left or right passage. Then the verifier randomly calls out a direction: ``Come out from the left!'' or ``Come out from the right!''

If the prover truly knows the spell, then no matter which side she entered from, she can pass through the magic door and emerge from whichever direction the verifier specifies. If she does not know the spell, she cannot pass through the door and can only retrace her steps.

Each time, she has only a 50\% chance of guessing correctly.

\vspace{1em}

In other words, each round of verification is like flipping a coin.

Repeat this process 20 times. If the prover succeeds every time, the verifier can be confident she knows the spell---the probability of cheating has dropped to less than one in a million ($0.5^{20} \approx 0.000001$). And throughout the entire process, the verifier has never heard a single syllable of the spell.

\vspace{1em}

This is a sorcery of revelation without exposure. It is called a zero-knowledge proof.

Its essence is to \textbf{prove the existence of knowledge without revealing the knowledge itself.}

What we display is the shadow of a capability, not the capability itself.

The verifier watches the prover emerge from the correct direction again and again; he becomes certain of her capability, yet knows nothing whatsoever about its content.

\begin{myquote}
This cave analogy comes from a classic 1989 paper by the cryptographer Quisquater et al., \textit{How to Explain Zero-Knowledge Protocols to Your Children}.

It became a classic because it reveals, in the simplest possible way, a profound truth:

\textbf{Proving and disclosing are two entirely different things.}

\end{myquote}

\newpage
\subsection{The Inverted Shadow of Plato's Cave: A Shadow Is Not the Thing Itself, Yet It Can Protect the Original}

From Ali Baba's cave, let us walk toward another, far more ancient one.

In \textit{The Republic}, Plato painted his famous Allegory of the Cave: prisoners are chained inside a cavern, their backs to a fire, able to see only the shadows on the wall. They mistake the shadows for reality, until one of them breaks free, walks out of the cave, and beholds the true sun and the true forms of things.

\vspace{1em}

Plato's teaching is that shadows are false; only the forms are real.

Those who seek truth should cast off the shackles of shadow and face the original directly.

\vspace{1em}

Yet in the Fifth Dimension, zero-knowledge proofs offer a radically different perspective:

In certain contexts, the shadow is all we need.

We do not need to hear the spell itself (the original).

We need only see the prover's ability to pass through the magic door (the shadow).

\vspace{1em}

This shadow already contains all the information we require---and it is safer and more useful than the original.
If the verifier were to hear the spell, he would possess the same power as the prover---and that would be dangerous for the prover.

\vspace{1em}

The shadow conveys the information while protecting the secret.

If we can complete every verification through shadows on the wall, why risk gazing at the original?

Staring directly at the sun scorches the eyes; staring directly at private data causes leaks.

\vspace{1em}

The shadow is a safer vessel for truth.

Sometimes, the shadow is more useful than the original.

Because it protects the original while transmitting the original's truth.

\begin{importantbox}
   The less you know, the less your liability, the lower the risk, the safer the collaboration.
\end{importantbox}

The Fifth Dimension embodies a survival wisdom: \textbf{from omniscience to the minimum necessary.}

This is not a compromise with truth, but the principle of minimizing truth's vessel.

\newpage
\subsection{The Separation of Knowledge and Content: Truth Can Be Proven, Yet Cannot Be Spoken}

When we say ``I know,'' we usually conflate two different things:

\textbf{Knowing the truth value}: I know the prover knows the spell---this is a judgment about truth.

\textbf{Knowing the content}: I know what the spell is---this is an acquisition of content.

In everyday experience, these two seem inseparable.

\begin{importantbox}
It is hard to believe something without understanding its content.
\end{importantbox}

Through zero-knowledge proofs, we can design systems that allow certain kinds of knowledge to flow (truth-value judgments) while blocking other kinds (content acquisition).

\vspace{1em}

Laozi wrote, ``The Dao that can be spoken is not the eternal Dao''---the way that can be put into words is not the enduring Way.

This line is usually read as a lament over the limits of language,

but in the context of zero-knowledge proofs, it acquires a new meaning:

\textbf{Some truths can be proven, yet cannot be spoken.}

\vspace{1em}

These may sound like philosophical abstractions, but they are in fact reshaping the trust architecture of the digital world through the hard craft of cryptography. With these techniques, we can let data be used and set in motion without revealing its details to those who use it.

We can prove to a bar that we are over 18 without showing our name, address, or date of birth on our ID; prove to a landlord that our bank balance exceeds 100,000 yuan without displaying the exact figure or account information; prove to an insurance company that we do not carry a certain hereditary condition without exposing our entire genome; and a corporation can prove to regulators that its operations are compliant without revealing trade secrets.

\vspace{1em}

What all these scenarios share is this: the verifier needs truth-value knowledge, not content knowledge. Cryptographers have used elliptic curves, polynomial commitments, recursive proofs, and other tools to transform this intuition of revelation-without-exposure into running code.

\vspace{1em}

\textbf{The First Form of the Stealth Codex, through precise separation, makes data usable without being leaked.}

\newpage

\section{Second Form of the Stealth Codex: The Art of Computation}

If the Art of Verification solves the problem of ``how to prove without exposing,'' then the Art of Computation tackles a harder challenge:
\textbf{how to collaborate without surrendering secrets in a world of mutual distrust.}
Once real work begins, privacy leaks along three paths.

\vspace{1em}

\textbf{Leaked to whom}: will the people involved in the computation see my raw data directly?

\textbf{What is leaked}: even if they cannot see the plaintext, can they still infer the content from the computation process?

\textbf{Exposed intent}: even if the content is secure, does the pattern of what I accessed, what I requested, what I chose already betray what I truly care about?

\vspace{1em}

The three layers of the Art of Computation's protective barrier are built along precisely these three fault lines:

First, hide the cards each party holds.
Then, hide the content that emerges during computation.
Finally, shroud even the act of attention and choice itself in mist.
This is not merely an upgrade in algorithms---
it is the deepening of the privacy boundary from ``who I am'' to ``what I want to do.''


\subsection{Leaked to Whom: Parties Who Distrust Each Other, Yet Compute Together}

Collaboration is society's most fundamental need. How to allow multiple parties to complete a joint computation while each guards their own secrets---that is the first problem collaboration must solve.

\phantomsection
\addcontentsline{toc}{manualsubsubsection}{The Millionaire's Problem: Not Knowing the Amounts, Yet Knowing Who Is Richer}
\subsubsection{The Millionaire's Problem: Not Telling the Other How Much I Have, Yet Knowing Who Is Richer}

In 1982, the Chinese-American scientist Andrew Yao, who would later receive the Turing Award, posed the famous Millionaire's Problem.
Two millionaires meet at a party.
Driven by a subtle vanity, each wants to know who is wealthier,
but neither is willing to reveal their exact fortune.

\vspace{1em}

This seemingly impossible problem---how can mathematics resolve it?

In the physical world, comparison requires disclosure.

To know who is heavier, both must step onto a scale whose reading is visible to all; to know who is richer, both must show their bank balances. Comparison presupposes the disclosure of information. In the Fifth Dimension, comparison can proceed through blind computation.

How does it work in practice?

A \textbf{simp}lified version of the idea: the two millionaires each encrypt their wealth figures and mix them together in a special way. This mixing guarantees that a comparison result (who is larger) can be derived, yet neither party's original wealth figure can be recovered.
The two millionaires obtain the conclusion ``A is wealthier than B,'' yet they will never know each other's exact fortune.

\vspace{1em}

They learned the result of the comparison without learning the content being compared.
The problem is elegantly resolved. The solution to this problem laid the technical foundation for secure multi-party computation.

The reason this matters is not the vanity of millionaires, but because it was the first clear demonstration that: \textbf{results can be shared without sharing the cards.}

\phantomsection
\addcontentsline{toc}{manualsubsubsection}{Secret Sharing: Giving You Part of the Treasure Map, Yet Revealing Nothing}
\subsubsection{Secret Sharing: Giving You Part of the Treasure Map, Yet Revealing Nothing}

Often, we still need to entrust a secret to others---the question is, how to hand it over without placing our fate in any single person's hands?

A pirate captain wants to entrust his treasure map to five lieutenants. He trusts no one among them---any one of them might steal the treasure if they held it alone---yet he cannot destroy the map either.
He needs a mechanism:
any five lieutenants working together can reconstruct the treasure map;
but any four conspiring together cannot glean even a fragment.

\vspace{1em}

In 1979, Adi Shamir produced an elegant solution: \textbf{secret sharing}.

The mathematical principle is this: it takes 2 points to determine a straight line, 3 points to determine a parabola, and $k$ points to determine a polynomial of degree ($k-1$).
Shamir's scheme hides the secret $S$ in the parameters of a polynomial curve, then randomly selects $n$ points on the curve and distributes them to $n$ people. Any $k$ people pooling their points can reconstruct the entire curve and recover the secret $S$. Fewer than $k$ people, however, find their points to be nothing but random noise---they do not obtain a partial secret; they have \textbf{not even a single clue}.

This all-or-nothing property means the secret cannot be glimpsed by piecing together fragments. What it rewrites is not the treasure map, but trust itself:
it does not slice the secret into pieces, but transforms ``the power to obtain the secret'' into something that can only happen through collaboration.

\begin{importantbox}
Trust need not be placed in any single point; it is distributed across the mechanism of collaboration.
\end{importantbox}



\newpage
\subsection{What Is Leaked: Homomorphic Encryption Lets Data Stay Hidden Yet Still Be Computed Upon}

Even when multi-party collaboration is feasible, we still encounter a second problem:
sometimes we are not computing jointly with another party, but handing our data to a more powerful third party for processing.
Here, the risk is no longer that participants peek at one another's cards, but that the content itself must be protected---computation must proceed in the dark.

\vspace{1em}

Imagine locking your X-ray inside a completely opaque black box and shipping it to a distant specialist.
Built into the walls of the box is a pair of gloves that reach inside.
The specialist cannot see what is in the box, but he can operate through the gloves.
Following a set procedure, he crops, rearranges, and applies chemical treatments to the film in total darkness---every operation is performed inside the sealed box.


At last he tells us the procedure is complete.
We receive the box and open it to find the X-ray has been processed with a diagnosis attached---and the specialist never actually ``saw'' our film.

\vspace{1em}

This is the sorcery the Stealth Codex calls \textbf{homomorphic encryption}:

\textbf{Computation is performed directly on encrypted data, and the decrypted result is identical to what plaintext computation would have produced.}

Data remains in ciphertext the entire time, yet it can be meaningfully processed---this is the ideal realization of usable yet invisible.
Mathematically, homomorphic encryption exploits a special algebraic structure: if an encryption scheme allows us to perform addition and multiplication on ciphertexts, and the decrypted results are equivalent to performing those operations on the plaintext, then the scheme is \textbf{homomorphic}.

\vspace{1em}

In 2009, Craig Gentry constructed the first fully homomorphic encryption\footnote{Craig Gentry, \textit{A Fully Homomorphic Encryption Scheme}, Stanford University, 2009.} scheme, conquering the challenge of supporting unlimited additions and multiplications simultaneously, and proving that arbitrary computation can be performed on encrypted data.

This means cloud computing can process data without ever knowing what the data contains.

What is uploaded is ciphertext; what the cloud performs is blind computation.

What is downloaded is an encrypted result that only we can decrypt.

\begin{importantbox}
The computation is completed, yet no information was ever seen by anyone else.
\end{importantbox}

\newpage
\subsection{Protecting Intent: Saying Nothing, Yet Already Revealing What You Were Looking At}

We often overlook the most intimate layer of all: the focus of our attention is itself private.

\vspace{1em}

During World War II, there was a classic case: even when the enemy perfectly encrypted every telegram's content, Allied intelligence officers could still analyze the frequency and locations of transmissions to precisely infer the movement routes of enemy command posts.

\vspace{1em}

\textbf{Sometimes, what you looked at is more lethal than what you saw.}

If a CEO frequently checks the financial statements of an otherwise unremarkable small company,

then even if he says nothing, the market can smell an acquisition in the air;

if someone accesses their ex's social media feed at four in the morning,

then the ever-climbing access count alone betrays the struggle within.

\vspace{1em}

Since no amount of encryption can conceal the act of access itself,

the only way to break the pattern is to make every action meaningless.

To hide the one door you truly want to knock on, you cannot simply stand before that door alone---

\textbf{you must walk down the corridor and knock on every single one.}

In the eyes of an observer, you are merely wandering about without any discernible pattern.

\vspace{1em}

When truth masquerades as falsehood, falsehood becomes indistinguishable from truth---the real intent dissolves into a fog of data.

This technique is called \textbf{oblivious access}---the elusive footwork of the Stealth Codex.

\begin{importantbox}
\textbf{To hide a single leaf, we must fabricate an entire forest.}
\end{importantbox}

Many privacy protocols in the real world are merely variations on this theme.

Some are used to confirm ``whether we share a common secret''; others ensure ``I've given you two choices, and you can take only one.'' Let us see how ingeniously these algorithms are designed.

% Casual readers need only remember this: \textbf{Intent itself can leak, and protecting intent often requires manufacturing additional fog.}
\newpage
\phantomsection
\addcontentsline{toc}{manualsubsubsection}{Private Set Intersection: No Exchanging of Lists, Yet Confirming What Both Sides Share}

\subsubsection{Private Set Intersection: No Exchanging of Lists, Yet Confirming What Both Sides Share}

If oblivious access is one person's self-preservation in the dark, then when two people meet in the dark forest, things become far more complicated.

A crime boss wants to cross-reference lists to root out a mole. The dilemma: without exchanging data, no connection can be found; but exchange the data directly, and both sides expose their hand. It is like the embrace of hedgehogs---longing to confirm what they share in common, yet afraid of being pierced by the other's spines.

\vspace{1em}

To illustrate the underlying principle, let us use the elegant dance of double-locking.

Imagine Alice and Bob each hold a unique private key (Lock A and Lock B).

Round One (Alice initiates): Alice places each piece of secret data into a box, fastens her Lock A, and sends the boxes to Bob.

\vspace{1em}

Round Two (Bob adds his lock and sends back): Bob does two things:

1. He adds his own Lock B to each box Alice sent, then sends them back to Alice.

\textbf{2. (The crucial step)} He places his own secret data into new boxes, locks them with his Lock B, and also sends them to Alice.

\vspace{1em}

Round Three (The reveal): Alice receives two sets of boxes.

She performs one action: she adds her Lock A to the boxes that carry only one lock (Lock B).

At this point, Alice has two sets---one derived from her original data, the other from Bob's data. Though they come from different sources, both are now secured by both Lock A and Lock B.

Mathematical commutativity guarantees this: locking with A first then B produces the exact same ciphertext fingerprint as locking with B first then A.

\vspace{1em}

Alice need not open a single lock. She simply compares the external appearance of the boxes. If two look identical, the data inside must be the same---that is the intersection. If they differ, she can never unlock them, and will never know what Bob's boxes contained.


\begin{importantbox}
Without either side knowing what the other holds, \textbf{both confirm they share a common secret.}
\end{importantbox}

\newpage
\phantomsection
\addcontentsline{toc}{manualsubsubsection}{Oblivious Transfer: Both Are Given to You, Yet You Can Open Only One}
\subsubsection{Oblivious Transfer: Both Are Given to You, Yet You Can Open Only One}

\textbf{Oblivious transfer} is the most fundamental building block of the Art of Computation---the starting point of all the magic.

Picture the iconic moment in \textit{The Matrix}: Morpheus extends both hands---in one, the red pill (the brutal truth); in the other, the blue pill (blissful illusion).

\vspace{1em}

Neo wants to choose one pill, but must never let Morpheus know which one he chose (protecting intent). At the same time, Morpheus is willing to offer the pills but will never allow Neo to take both (protecting assets). In the physical world, this seems impossible---Morpheus must watch Neo take one pill to ensure he did not take the other.

\vspace{1em}

To solve this riddle, they devise an ingenious protocol:

Round One (Morpheus initiates): Morpheus hands Neo a special piece of metal and says, ``This metal will automatically generate two locks---one for the red box, one for the blue box. It has only enough material to forge a single key, but you get to decide which lock it fits.''

Round Two (Neo's binary choice): Neo wants the red pill, so he must fashion the only key for the red box's lock.

Round Three (Morpheus operates blind): Neo returns both locks to Morpheus. Morpheus places the red pill in the red box, the blue pill in the blue box, then pushes both locked boxes to Neo.

Round Four (The reveal): Neo takes out the only key, opens the red box, and obtains what he wanted. As for the blue box---he can never open it.

This is the deepest primitive of secure computation---the essence of oblivious transfer:

\textbf{I give you two locks, but only one key.}

\begin{importantbox}
I gave you everything, yet I am certain you can take only half;

you got what you wanted, but I have no idea which half you took.

\end{importantbox}

\newpage
% ------------------------------------------------------------
\section{Third Form of the Stealth Codex: The Art of the Macro}
% ------------------------------------------------------------

The Art of Verification protects the input---who you are;
the Art of Computation protects the process---what you used;
we still need to guard---the output.
When we extract statistical patterns from data and generate analytical reports, how do we ensure the results do not leak back to expose individual privacy?

This is precisely the data-security barrier the Art of the Macro seeks to construct.

\subsection{Differential Privacy: Given Everyone's Data, You Can See Only the Whole, Never the Individual}

In the Second Form, we learned: to hide a single leaf, you must fabricate an entire forest.

That was \textbf{hiding the part within the whole}---drowning real intent in a sea of decoys.

The Third Form inverts the logic: to protect every single tree, we \textbf{allow only the forest to be seen}.

This is the core technique the Art of the Macro employs---the essence of \textbf{differential privacy}\footnote{Cynthia Dwork, \textit{Differential Privacy}, ICALP, 2006.}.

\vspace{1em}

In the real world, organizations use this principle to collect statistical signals. The method borders on vandalism: a small amount of random noise is added to each person's true answer. If the true answer is ``yes,'' the system may record ``no''; if the true age is 35, the system may record 38.

Every tree is deliberately shaken, blurred beyond recognition; noise drowns out signal.

\vspace{1em}

\textbf{The miracle happens when we step back and survey the entire forest.}

When a million data points converge, the noise cancels itself out, and statistical patterns emerge naturally.

Random noise is symmetric---some trees sway left, some sway right,

and under idealized design, their average deviation tends toward zero.

The shaking of individuals vanishes at the macro level, while the contour of the forest becomes crisp at the macro level.

\begin{importantbox}
Hide a leaf with the forest---\textbf{let the part disappear into the whole.}

Protect each tree with the forest---\textbf{let the whole be observed in place of the part.}
\end{importantbox}

The Third Form and the Second Form arrive at the same destination by different paths: both use the many to protect the one, trading precision for safety.


% ============================================================
\section{The Price: Fuel for the Alchemist's Furnace}
% ============================================================

Having studied all three forms of the Stealth Codex, we might ask:

If these techniques are so ingenious and were invented years ago, why have they not yet achieved widespread adoption in the Fifth Dimension?


The answer lies not only in algorithmic elegance, but in the \textbf{cost of energy}.

\vspace{1em}

Computing on ciphertext often costs orders of magnitude more than computing on plaintext.

This is first and foremost a mathematical and engineering expense, not some free spell.

The universe's natural tendency is entropy increase: information tends to leak, to diffuse, to blur.

To keep privacy from leaking, enormous energy must be continuously poured in.

\vspace{1em}

Medieval alchemists believed that with enough fire, lead could be transmuted into gold.

They were wrong---but their intuition was half right: \textbf{transformation requires energy.}

\vspace{1em}

Privacy-preserving computation is the alchemy of the Fifth Dimension.

We cast raw, exposed data into the furnace and temper it in the roaring flames of computational power, until it becomes the encrypted gold that is usable yet invisible.

\vspace{1em}

\textbf{The hotter the furnace, the purer the gold. The higher the cost, the truer the value.}

Precisely because it is expensive, it filters out the data that is truly worth protecting.

Only the information for which someone is willing to burn computational power---that is the real gold mine.

\begin{thinkbox}[Reflection]
When we wield these costly arts, what changes will they bring to society?
\end{thinkbox}

\newpage
% ============================================================
\section{Order: Rebuilding Trust in a Distributed World}
% ============================================================
The reshaping of society's trust structure is the first transformation the Stealth Codex brings.

\subsection{The Byzantine Generals' Siege}

Let us return to the year 1453. The armies of the Ottoman Empire lay siege to Constantinople, with multiple contingents encamped outside the city walls. The generals need to coordinate their actions---either all attack or all retreat; any inconsistency will spell disaster.

The problem: the generals can communicate only through messengers, and messengers may be intercepted or may defect. Worse still, some of the generals themselves may be traitors, deliberately sending contradictory messages to sow chaos.

\vspace{1em}

In such an environment saturated with distrust, how can consensus be reached?

This is the Byzantine Generals Problem, posed by the computer scientist Leslie Lamport and collaborators in 1982\footnote{Leslie Lamport, Robert Shostak, Marshall Pease, \textit{The Byzantine Generals Problem}, ACM TOPLAS, 1982.}. It is not a historical question but a fundamental question about distributed systems. Lamport proved a theorem: when the number of traitors exceeds one-third of the total, \textbf{no algorithm can guarantee consensus.}
This one-third boundary is hard, mathematical, and impassable.

Conversely, this also means: as long as honest nodes hold a sufficient majority, and a suitable protocol is in place, consensus can be reached without any central authority.

\vspace{1em}

Hobbes, in \textit{Leviathan}, depicted the state of nature: without government, every man is an enemy to every man, and life is ``solitary, poor, nasty, brutish, and short.'' To escape this condition, people surrender their power to an absolute sovereign---the Leviathan---who maintains order.

The Leviathan is a centralized solution: a powerful central node that resolves the trust problem.

Byzantine consensus algorithms offer another possibility: a distributed solution.

\vspace{1em}

Trust can emerge from the structure of the network, guaranteed by protocol rules.

This is not a rejection of centralization---in many cases, centralization is efficient, necessary, even the superior choice. But the Stealth Codex grants us \textbf{the power to choose}: we have the ability to build a reliable distributed trust system without any central node.

\newpage
\subsection{Cold Trust and Warm Trust}

We often say that blockchains and privacy-preserving computation are machines that manufacture trust.

This may be a beautiful misunderstanding.

More precisely, they are machines that minimize the need for trust.

Through years of research in privacy-preserving computation, I have come to realize that what we casually call ``trust'' actually encompasses two very different things. Two words are often used interchangeably: \textbf{trust} and \textbf{confidence}.

\vspace{1em}

\textbf{Trust} is rooted in human relationships; it contains vulnerability.

I believe my neighbor will not rob me even if I leave the door unlocked---this is a gamble, a leap of faith.

It is precisely because I can be hurt that my trust carries moral weight.

\vspace{1em}

\textbf{Confidence} is rooted in systemic reliability.

I leave the door unlocked because I have an automated defense system.

I depend not on my neighbor's morality but on the certainty of the system.

When betrayal becomes mathematically impossible (or its cost approaches infinity), trust effectively vanishes, replaced by \textbf{a confidence level}.

\vspace{1em}

Let us distinguish two kinds of trust:

\textbf{Warm trust}---trust between human beings. It has warmth, fragility, the possibility of betrayal. It is precisely because betrayal is possible that forgiveness has meaning; precisely because there is fragility that relationships can deepen; precisely because there is risk that trust is precious;

\textbf{Cold trust}---trust between a person and a protocol. It is solid, reliable, verifiable, yet it has no warmth, no elasticity, no possibility of redemption. We cannot forgive a smart contract; we cannot reconcile with a cryptographic protocol.

\begin{importantbox}
Mathematics has eliminated the \textbf{floor of betrayal}, and in doing so, has locked down the \textbf{ceiling of goodwill}.
\end{importantbox}

What we gain in the Fifth Dimension is not a utopia brimming with trust, but a mechanical paradise that runs with flawless precision---without needing trust at all.

\newpage
\subsection{The Credit of the Nameless}
For thousands of years, the trust systems of human civilization have been anchored to identity.

\vspace{1em}

Trust is a costly form of social capital.

We trust someone often because he is somebody---bound by family ties, endorsed by institutions, weighted by the sediment of history. His past serves as collateral for his future.
This kind of trust is \textbf{deeply personalized}, dependent on the accumulation of reputation, heavy and difficult to transfer.

\vspace{1em}

The Fifth Dimension brings credit without identity.
In the world built by the Stealth Codex, trust no longer needs to know us---it only needs to verify us.
Whether we are a multinational giant or an anonymous hacker, so long as the mathematics checks out, the credit is complete.

\begin{importantbox}
    \textbf{Before computational power, all are equal.}
\end{importantbox}

This is a paradigm shift from permission-based trust to verification-based trust:

\textbf{Permission-based trust} is a soft constraint, grounded in law and morality.

Because the physical threshold for wrongdoing is low, it relies primarily on after-the-fact legal punishment to deter betrayal.

\textbf{Verification-based trust} is a hard constraint, grounded in mathematics and code.

Because the computational cost of wrongdoing approaches infinity, betrayal becomes physically infeasible.

\vspace{1em}

We have crossed the watershed:

from relying on the other party's subjective goodwill (believing they will not act in bad faith),

to relying on objective mathematical certainty (verifying that they cannot act in bad faith).


\newpage

\section{Value: A New Paradigm of Property---From Owning to Using}

In the physical world, land, gold, oil, and other resources obey the principles of conservation and exclusivity.

This is a zero-sum game---if I have it, you do not.

Millennia of evolution have taught us: security comes from hoarding, power comes from walls.

\vspace{1em}

But in the Fifth Dimension, the Stealth Codex has triggered a \textbf{phase transition toward liquidity}.

Applying the deed-and-title logic of land management to define property rights in this new realm is the folly of marking the side of a moving boat to find a dropped sword.

\begin{thinkbox}[Reflection]
    Since copying is virtually free, what does ownership even mean?

    If value comes from flow, how should property rights be defined?
\end{thinkbox}

\subsection{The Music Industry's Lesson: Flow Wins}

Looking back at 20 years of music copyright evolution, what we see is a \textbf{war between ownership and access}.

\textbf{The solid-state logic of the physical world:} Record labels sold physical media (vinyl, CDs); users paid for exclusive copies. Security meant a shelf full of albums.


\textbf{The liquid-state logic of the Fifth Dimension:} MP3 drove the cost of copying to zero; Napster tore down the walls; and finally Spotify took the stage---it does not own the copyrights to songs, yet it commands \textbf{connection} to the entire world.

\vspace{1em}

\textbf{The winner is decided!}

The once-invincible record-industry titans no longer monopolize the gateway. Streaming platforms that control user touchpoints through connection, recommendation, and distribution have become new centers of power.

In a liquefied society, possessing an MP3 file is dead data. Only when that song is surfaced by algorithms, heard by millions, added to countless playlists does it come alive.

\vspace{1em}

\textit{Value no longer springs from \textbf{static inventory}, but from \textbf{dynamic circulation}.}

\newpage

\subsection{The Smelter's Gravity}

This is not a comfortable new world.

The Stealth Codex solves the technical question of ``can we exchange,'' but for flow to actually happen, it must also resolve ``are we willing to exchange.'' This demands that we accept a counterintuitive premise:

\textbf{Letting go creates more value than holding on.}

\vspace{1em}

Traditional data companies imagine themselves as lords of a mountain of gold: I own the data, and if you want to use it, pay up.
We must acknowledge a reality: the vast majority of raw data, before it has been refined, is worth close to zero---or is even a liability. Hoarding piles of uncleaned, unlinked, context-free raw data is like filling a backyard with iron ore---it looks like wealth, but all you truly own is dust and occupied space.

\vspace{1em}
The better positioning is that of a \textbf{smelter}.

Public data, de-identified data, synthetic data---these are like ore scattered across the open plain, accessible to many. The real value lies not in who possesses the ore, but in who can forge it into refined steel.

\vspace{1em}

The value of data is not in pricing the stone, but in who owns the blast furnace.

If our algorithmic models can amplify the value of data a hundredfold, while data owners themselves can amplify it only once, then data will naturally flow toward us. Data owners provide not an expensive commodity but cheap raw material; and what they seek in return is not merely cash, but the refined product: more precise diagnoses, more efficient matching, wiser decisions.

\vspace{1em}

The one who walls things off sits atop a mountain of gold, besieged and helpless;

the one who opens up has no walls, yet data flows willingly into his furnace.

This is no longer a zero-sum game that creates nothing new, but an incremental exchange that amplifies value.

\begin{importantbox}
Flow creates value.
\end{importantbox}

Only when data flows---computed upon, iterated, imbued with meaning---does it truly become an asset.


\newpage
% ============================================================
\section{Existence: The Boundary Between Knowing and Not Knowing}
% ============================================================

\begin{myquote}
\textit{The Emperor of the Southern Sea was called Swift, the Emperor of the Northern Sea was called Sudden, and the Emperor of the Center was called Chaos.}

\textit{Swift and Sudden would from time to time meet in the land of Chaos, and Chaos treated them with great kindness. Swift and Sudden wished to repay the virtue of Chaos,}

\textit{and said: ``All people have seven openings through which they see, hear, eat, and breathe. Chaos alone has none. Let us try to bore some for him.''}

\textit{Each day they bored one opening. On the seventh day, Chaos died.}
\end{myquote}


Swift and Sudden stand for clarity and order;
Chaos stands for the blurred and the indistinguishable.

When we try to pierce chaos with clarity, chaos dies.

\vspace{1em}

The French mathematician Laplace was a disciple of Swift and Sudden.

He envisioned an omniscient demon---if it knew the position and momentum of every particle in the universe, it could predict all of the future.
Twentieth-century physics shattered this crystal ball---the world is inherently uncertain.
Some things, once seen clearly, cease to be themselves.

\vspace{1em}

Looking back at the three forms of the Stealth Codex, we discover their shared secret weapon---\textbf{manufactured chaos}.

\textbf{Differential privacy} is noise woven into a precise melody;

\textbf{Secret sharing} is a whole truth shattered into random fragments;

\textbf{Zero-knowledge proofs} are random challenges injected into an interactive labyrinth.

\vspace{1em}

Chaos does not mean the absence of boundaries. A self without boundaries dissolves; a self whose boundaries are repeatedly violated is traumatized. In the physical world, boundaries are tangible. Skin is the boundary of the body; walls are the boundary of private space. In the Fifth Dimension, these boundaries are gradually vanishing. Locations are tracked, spending is logged, social ties are analyzed, desires are inferred\ldots

% We can choose who enters our space, who glimpses our secrets.

% That power of choice itself is a part of the self.

% We become transparent---not because we chose transparency, but because we are powerless not to be.
\vspace{1em}

When boundaries vanish, the self begins to dissolve. We can no longer be sure which thoughts are our own and which were planted by algorithms; we can no longer be sure which choices are autonomous and which were engineered.

\textbf{Having lost the ability to say ``no,''} ``yes'' becomes meaningless.

\vspace{1em}

The significance of the Stealth Codex for the self lies precisely here:
it is not merely information-security technology, but an engineering project to rebuild the boundaries of the self in a world of data.
Boundaries do not mean isolation; they are the precondition for connection.

\subsection{Negative Capability: The Wisdom of Not Looking}
In 1817, the English poet Keats, writing to his brothers, coined a curious concept: Negative Capability.
He said that true genius is the capacity to dwell in uncertainty, mystery, and doubt, without any irritable reaching after fact and reason.

\vspace{1em}

Shakespeare possessed this capacity, which is why Hamlet can hang forever between action and hesitation,
why \textit{The Tempest} can shimmer on the boundary between magic and reality, and the audience feels no anxiety---because irresolution itself is the meaning.

\vspace{1em}

We live in an age of positive-capability surplus.

We can track every person's movements, predict every desire\ldots

Technology grants us an almost limitless power to know.

\textbf{But ``can'' has never been the same as ``should.''}

In an age of technological explosion, this has become a rare form of discipline.

\vspace{1em}

Zhuangzi wrote: ``My life has a limit, but knowledge has none. To pursue the limitless with the limited---that is perilous indeed.''

To chase infinite knowledge with a finite life is a doomed endeavor.

What the Stealth Codex teaches us is \textbf{the wisdom of restraint}---the ability to choose not to see.

\textbf{True wisdom is knowing what should be known, and what should remain unknown.}

\begin{thinkbox}[Reflection]
If you had the power to see into everyone's heart, would you use it?

And if you did, could you still have a true friend or a true lover?
\end{thinkbox}

\newpage
% ============================================================
\section{Coda: Love Is the Anti-Zero-Knowledge Proof}
% ============================================================

If, in the future, we quarrel with our beloved and they demand: ``Do you truly love me?''

We pull out our phone, run a zero-knowledge proof protocol, and prove that our heart-rate fluctuations, dopamine levels, and attention-allocation patterns perfectly match the statistical signature of love---while exposing no raw psychological data whatsoever.
They stare at the green ``Verification Passed'' on the screen.

In that moment, love is dead.

\vspace{1em}

Because the essence of love is precisely the \textbf{anti-zero-knowledge proof}.

Love is not proof; it is exposure. Love is the voluntary surrender of encryption, the willing embrace of vulnerability.

To love someone means we are willing to let them see:

the ugliest parts, the most fragile parts, the parts unfit for anyone's eyes.

I hand over my fragility; you let me see your darkness.

And then, from each other's broken pieces, we generate a world that belongs to us alone---

a singular, unrepeatable existence we have created together.

\vspace{1em}

And so the zero-knowledge-proof version of ``I love you'' is not love.

It proves we possess the capacity for love, but precisely because nothing was exposed, no actual loving took place.

\textbf{It is because it cannot be verified, yet we choose to believe anyway, that this gamble constitutes the weight of trust.}

\vspace{1em}

The Stealth Codex teaches us to put on armor, but it cannot teach us to embrace the ones closest to us.

We practice the art of invisibility to survive in a dangerous world,

\textbf{but in the end, we will take off the armor.}

So that before that one particular person, we may appear---whole, unreserved, and utterly visible.

\textbf{``Know the white, yet keep to the black.''}

\begin{importantbox}
To guard the darkness is to preserve our right to choose---

when, where, and before whom to reveal that beam of light called the self.
\end{importantbox}
